%%%	Vorlage für Bachelor-/Masterarbeit
%%%
%%%	Autor: 		Nicolas Peschke
%%%	Version: 	2.2
%%%	Datum:		09.07.2018
%%%
%%%	Die Vorlage kann unter Studenten der Molekularen Biotechnologie der
%%%	Universität Heidelberg weitergegeben und verwendet werden.
%%%
%%%	Für die Richtigkeit und Konformität des Layouts mit den Richtlinien
%%%	des IPMBs wird keine Haftung übernommen, BITTE SPRECHT AUF JEDEN
%%%	FALL DAS LAYOUT EURER ARBEIT MIT EUREN BETREUERN AB!!!
%%%
%%%	Um die Titelseite und die zweite Seite richtig zu setzen müsst ihr
%%%	nur in dieser Datei unter "Definitionen für Titelseite" die auf
%%%	euch passenden Eintragungen vornehmen.
%%%	In der Datei declaration.tex muss dann noch die Anschrift der Prüfer
%%%	richtig angegeben werden.
%%%
%%%	Für das Zitieren muss im Befehl
%%%	\addbibresource{bib/bachelorarbeit.bib} durch den
%%%	Dateinamen und Dateipfad eurer BibTeX Datei aus Citavi an.

\documentclass[parskip, twoside]{scrreprt}			% Dokumentenklasse scrreprt, für beidseitigen Druck, deutsche Absatzabstände

%%%%%%%%%%%%%%%%%%% Pakete %%%%%%%%%%%%%%%%%%%%%%%%%%%%%%%%%%%%%%%%%%%%%%%%

%%% Sprache und Eingabecodierung %%%
\usepackage[greek, ngerman, main=english]{babel}
% Sprachpaket greek damit griechische Buchstaben (z.B. im Literaturverzeichnis) keine Probleme machen mit main= die standardsprache für das Dokument festlegen, damit werden automatische Überschriften wie Inhaltsverzeichnis/Table of Contents in der richtigen Sprache gesetzt

\usepackage[utf8]{inputenc}							% codierung damit man Umlaute direkt schreiben kann
\usepackage[T1]{fontenc}							% auch was mit codierung zu tun

%%% Schriftart und Form %%%
\usepackage{lmodern}								% Schriftart LatinModern Vektorisiert für schöneres PDF
\usepackage[onehalfspacing]{setspace}				% Korrekter Zeilenabstand nach IPMB Vorgabe
\usepackage[left=2.50cm, right=2.50cm, top=2.50cm, bottom=2.50cm, bindingoffset=10mm, includehead, includefoot]{geometry}
% Seitenränder nach Vorgabe von IPMB, mit Bindezugabe von 10mm, Kopf- und Fußzeile im Textfeld anstatt im Rand => 2,5cm zu Kopfzeile anstatt zum Text
\usepackage[headsepline]{scrlayer-scrpage} 			% schönere Kopfzeile
\usepackage{url}

%%% Bibliographie mit biber %%%
\usepackage[backend=biber, style=authoryear, giveninits=true, maxbibnames=99, uniquename=init, maxcitenames=2, hyperref=true, date=year]{biblatex}			% Benutze Biber/BibLaTeX zum Zitieren
\usepackage{xpatch}
\usepackage{csquotes}
\addbibresource{bib/quellen.bib}					% Pfad zur BibTeX Datei aus Citavi
% HIER EINS VON BEIDEN AUSWÄHLEN
\usepackage{bib/cellstyle}								% Cell Stil
%\usepackage{bib/cellstyle_it}							% Cell Stil mit kursivem et. al.

%%% Chemie, Biologie, Mathematik %%%
\usepackage{dnaseq}									% Für einzelne Sequenzen
\usepackage{texshade}								% Für alignierte Sequenzen
%%%%%%%%%%%%%%%%%% !!! ACHTUNG !!! dnaseq vor texshade vor mhchem vor graphicx einbinden !!! ACHTUNG !!! %%%%%%%%%%%%%%%%%%

\usepackage[version=4]{mhchem}						% Chemische Formeln
\usepackage{siunitx}								% SI Einheiten
\usepackage{amsmath}								% für schönere Formeln
\usepackage{listings}								% um Code zu setzen

%%% Tabellen und Abbildungen %%%
\usepackage{booktabs} 								% für schöne Linien in Tabellen
\usepackage{longtable}								% Tabelle mit Seitenumbrüchen
\usepackage{multirow}								% Mehrere Reihen in einer Tabelle zusammenfassen
\usepackage{rotating}								% für Tabellen im Querformat
\usepackage{subfigure}								% für mehrere Unterabbildungen in einer Abbildung
\usepackage{graphicx}								% für Größenänderungen für Grafiken
\usepackage{float}									% um Floats (z.B. Bilder) an Ort und Stelle mit [H] zu fixieren


\usepackage{acronym}								% Abkürzungen

\usepackage{lipsum}									% Blindtext zum Testen

\usepackage[hidelinks]{hyperref}
% Hyperlinks [hidelinks] um linkboxen zu deaktivieren (möglichst zum Schluss einbinden, da es sich anderen Paketen anpasst)

%%%%%%%%%%%%%%%%%%% KOMA-Script Einstellungen %%%%%%%%%%%%%%%%%%%%%%%%%%%%%
\usepackage{scrhack}								% Anpassung von anderen Paketen an KOMA-Script zur besseren Kompatibilität. Zuletzt Einbinden!

\emergencystretch=50pt								% Vergrößert Wortzwischenraum notfalls, wenn sonst die Zeile überläuft

\clubpenalty = 10000								% Verhindern von Schusterjungen
\widowpenalty = 10000								% Hurenkindern
\displaywidowpenalty = 10000						% und vor Formeln

% Einstellen der Kopfzeile
\automark[section]{chapter}							% automatisches Setzen von Kollumnentiteln in der Kopfzeile (chapter links, section rechts)
\renewcommand*{\chaptermarkformat}{} 				% Chapter ohne Nummern
\renewcommand*{\sectionmarkformat}{}				% Sections ohne Nummern

%%%%%%%%%%%%%%%%%%% Schriftart für Überschriften einstellen %%%%%%%%%%%%%%%
\setkomafont{title}{\sffamily}						% Festsetzen der Überschriftart auf Sans Serif
\setkomafont{disposition}{\usekomafont{title}}		% Festsetzen der weiteren Unterüberschriften auf Sans Serif
\setkomafont{author}{\usekomafont{title}}
\setkomafont{date}{\usekomafont{title}}
\setkomafont{caption}{\sffamily\small}				% Festsetzen der Captions auf Sans Serif klein
\setkomafont{captionlabel}{\usekomafont{caption}\bfseries\small}	% Festsetzen der Captionlabels auf Sans Serif Fett klein
\setkomafont{pagehead}{\normalfont\scshape}			% Festsetzen der Kopfzeilen auf normale Kapitälchen

%%%%%%%%%%%%%%%%%%% Definitionen für Titelseite %%%%%%%%%%%%%%%%%%%%%%%%%%%
%%% Hier müssen alle Anpassungen für euch vorgenommen werden,			%%%
%%% diese werden dann automatisch im Dokument angewendet				%%%
%%%%%%%%%%%%%%%%%%%%%%%%%%%%%%%%%%%%%%%%%%%%%%%%%%%%%%%%%%%%%%%%%%%%%%%%%%%
\newcommand{\workauthor}{NAME}
\newcommand{\worktitle}{TITEL}
\newcommand{\worksubtitle}{ZWEITETITELZEILE}		% Für mehrzeilige Titel, leerlassen wenn nicht benötigt
\newcommand{\worksubsubtitle}{DRITTETITELZEILE}
\newcommand{\studentid}{MATRIKELNUMMER}
\newcommand{\birthplace}{GEBURTSORT}
\newcommand{\workplace}{ARBEITSORT}
\newcommand{\workstart}{BEGINNARBEIT}
\newcommand{\workstop}{ENDEARBEIT}
\newcommand{\submissiondate}{ABGABEMONATJAHR}
\newcommand{\institute}{INSTITUT}
\newcommand{\department}{ABTEILUNG}
\newcommand{\degree}{BACHELOR/MASTER}				% Hier entweder "Bachelor" oder "Master" je nach Art der Arbeit eintragen
\newcommand{\worktype}{\degree{}arbeit}				% Hier braucht nichts verändert zu werden

%%%%%%%%%%%%%%%%%%% Paketeinstellungen %%%%%%%%%%%%%%%%%%%%%%%%%%%%%%%%%%%%%

\lstset{							% Einstellungen für listings
	frame=tb,						% Horizontale Linien oben und unten
	numbers=left,					% Zeilennummern auf der linken Seite
	basicstyle=\ttfamily\small,		% Schrift Monospace klein
%	language=,						% Festlegen der Programmiersprache
	tabsize=2,						% Tabs sind 2 Leerzeichen breit
}

\sisetup{							% Einstellungen für SI-Einheiten
	per-mode=fraction,				% macht Brüche bei Einheiten
	output-decimal-marker={,},		% gibt "," als Dezimaltrennzeichen aus
	separate-uncertainty,			% Unsicherheit mit +-
	table-align-uncertainty = true,	% Ausrichtung am +- in Tabellen
}

\begin{document}

%%%%%%%%%%%%%%%%%%% Titlepage für Bachelor/Masterarbeit %%%%%%%%%%%%%%%%%%%
\begin{titlepage}
	\centering
	{\Large Ruprecht-Karls-Universität Heidelberg\\
		Fakultät für Biowissenschaften\\
		\degree{}studiengang Molekulare Biotechnologie\\}

	{\vspace{\stretch{2}}}
	{\usekomafont{title}

		{\Huge \worktitle{}}

		{\Huge \worksubtitle{}}

		{\Huge \worksubsubtitle{}}

	}

	\vspace{\stretch{2}}
	{\Large \worktype{}}

	\vspace{\stretch{2}}

	{\Large
		\begin{tabular}{rl}
			Autor &\workauthor{}\\
			Geburtsort &\birthplace{}\\
			Abgabetermin &\submissiondate{}\\
		\end{tabular}
	}

	\vspace{\stretch{1}}

\end{titlepage}
%%%%%%%%%%%%%%%%%%% /Titlepage für Bachelor/Masterarbeit %%%%%%%%%%%%%%%%%%
\pagestyle{plain.scrheadings}		% Text bis zur Einleitung ohne Kopfzeile

\cleardoublepage
\thispagestyle{plain}
Die vorliegende \worktype{}  wurde im \institute{} in der Abteilung \department{} in der Zeit von \workstart{} bis \workstop{} angefertigt.

\begin{center}												% Hier muss die Anschrift für den/die Gutachter eingefügt werden
	\begin{tabular}{rl}
		Gutachter der Arbeit 	&Dr. Matthias Schlesner\\
								&Computational Oncology\\
								&eilslabs, Deutsches Krebsforschungszentrum\\
								&\\ % Nur für Master
						%		Zweitgutachter der Arbeit &Dr. Matthias Schlesner\\
						%		&Computational Oncology\\
						%		&eilslabs, Deutsches Krebsforschungszentrum\\
	\end{tabular}
\end{center}










\vspace{\stretch{4}}
Ich erkläre hiermit ehrenwörtlich, dass:
\begin{enumerate}
	\item ich die vorliegende \worktype{} selbständig unter Anleitung verfasst und keine anderen als die angegebenen Quellen und Hilfsmittel benutzt habe
	\item die Übernahme wörtlicher Zitate aus der Literatur und dem Internet sowie die Verwendung der Gedanken anderer Autoren an den entsprechenden Stellen innerhalb der Arbeit gekennzeichnet wurde
	\item ich meine \worktype{} bei keiner anderen Prüfung vorgelegt habe.
\end{enumerate}

Ich bin mir bewusst, dass eine falsche Erklärung rechtliche Folgen haben wird. 

% 3,5cm Abstand nach oben
\vspace{\stretch{1}}
% Ort, Datum
\workplace{}, der \workstop{}
\vspace{1,5cm}
% 8 cm breite Linie für die Unterschrift
\begin{center}
	\begin{minipage}[t]{8cm}
		% gepunktete Linie
		\centering \dotfill \\
		% Text unter der Linie
		\workauthor{}
	\end{minipage}
\end{center}


\vspace{\stretch{1}}

%%% Danksagung (optional) %%%%%%%%%%%%%%%%%%%%%%%%
\renewcommand\abstractname{\Large Acknowledgments}
\begin{abstract}


\end{abstract}
%%%%%%%%%%%%%%%%%%%%%%%%%%%%%%%%%%%%%%%%%%%%%%%%%%

\renewcommand\abstractname{\Large Zusammenfassung}
\begin{abstract}								% Laut IPMB Richtlinien kommt der deutsche Abstract immer zuerst!
	\begin{otherlanguage}{ngerman}				% jenachdem in welcher Sprache die Arbeit verfasst ist, muss die Umgebung für die jeweils andere Sprache gesetzt werden
		Hier steht ein deutscher Abstract!		
	\end{otherlanguage}						
\end{abstract}

\renewcommand\abstractname{\Large Abstract}
\begin{abstract}
	This is the english abstract!

\end{abstract}

\tableofcontents

{\setlength{\parskip}{0.2cm}		% Kleinerer Zeilenabstand zwischen den Einträgen
%	Abkürzungsverzeichnis
\addchap{Abbreviations}
	\begin{acronym}[ABCDEFGHIJK]	% Länge des Strings in den [..] entspricht dem Abstand der zwischen Acronym und Vollname gelassen wird
									% => mindestens so lange wie das längste Acronym
		% A B C D E F G H I J K L M N O P Q R S T U V W X Y Z		
		% Abkürzungen
		%\acro{acronym}{full name}
		
		% Formelzeichen
		%\acro{acronym}[$ Formelzeichen $]{full name}
		
		% als benutzt markierte Acronyme	
		%\acused{acronym}
		
	\end{acronym}
}

\chapter{Introduction}
\pagestyle{scrheadings} 	% Ab hier mit Kopfzeile


\chapter{Material and Methods}
	
\chapter{Results}
	
\chapter{Discussion}
	
\chapter{Outlook}

\appendix
\chapter{title}


\printbibliography


\end{document}
